% Generated draft: requires mathematical and source review.
\documentclass[11pt]{article}
\usepackage[margin=25mm]{geometry}
\usepackage{fontspec}
\setmainfont{Latin Modern Roman}
\usepackage{amsmath,amssymb,mathtools,unicode-math}
\setmathfont{Latin Modern Math}
\usepackage{enumitem,needspace,xurl,hyperref}
\setlength{\parindent}{0pt}
\setlength{\parskip}{5pt}
\setlength{\emergencystretch}{3em}
\newcommand{\N}{\mathbb{N}}
\newcommand{\Z}{\mathbb{Z}}
\newcommand{\Q}{\mathbb{Q}}
\newcommand{\R}{\mathbb{R}}
\newcommand{\C}{\mathbb{C}}
\newcommand{\F}{\mathbb{F}}
\newcommand{\A}{\mathbb{A}}
\newcommand{\PP}{\mathbb{P}}
\newcommand{\E}{\mathbb{E}}
\newcommand{\eps}{\varepsilon}
\newcommand{\dd}{\,\mathrm{d}}
\begin{document}
\section*{UnsolvedMath ID: 700248}
\textbf{Source problem code:} catalogue entry
\subsection{Definitions and mathematical statement}
Let \(C\) be a fixed nontrivial group; the supplied example is \(C=\Z/2\Z\). Let \(A\) and \(B\) be arbitrary groups. Under the convention that a split extension of \(B\) by \(A\) has normal subgroup \(B\) and quotient \(A\), such an extension is a split exact sequence
\[
1\longrightarrow B\longrightarrow E\overset{p}{\longrightarrow}A\longrightarrow 1,
\qquad s\colon A\longrightarrow E,\quad p\circ s=\operatorname{id}_A.
\]
After choosing the splitting and identifying the kernel with \(B\), this gives \(E\cong B\rtimes_\alpha A\), where \(\alpha\colon A\to\operatorname{Aut}(B)\) is a homomorphism and
\[
(b,a)(b',a')=(b\,\alpha(a)(b'),aa').
\]
The source does not explicitly fix its convention for the phrase ``extension of \(B\) by \(A\)''; the sequence above makes the convention used in this draft explicit.

The notation \(B\wr C\) denotes a wreath product, but the supplied record does not specify whether it is restricted or unrestricted. The regular restricted wreath product is \(B^{(C)}\rtimes C\), where \(B^{(C)}\) consists of functions \(C\to B\) with finite support. The regular unrestricted wreath product is \(B^C\rtimes C\), using all functions \(C\to B\). In both cases multiplication in the base group is pointwise and the shift action is \((c\cdot f)(x)=f(c^{-1}x)\). These possibilities can differ when \(C\) is infinite. Write \(W=B\wr C\) for the wreath product intended by the cited source, and let \(\operatorname{Aut}(W)\) be its group of automorphisms.

The supplied background attributes to Merzlyakov a construction embedding all split extensions, in a certain unified way, into \(A\times\operatorname{Aut}(W)\). Under the extension convention above, denote the embeddings from that construction symbolically by
\[
\iota_\alpha\colon B\rtimes_\alpha A\hookrightarrow A\times\operatorname{Aut}(W),
\qquad \alpha\in\operatorname{Hom}(A,\operatorname{Aut}(B)).
\]
The problem asks for a description of how the image subgroups
\[
\iota_\alpha(B\rtimes_\alpha A)\leq A\times\operatorname{Aut}(W)
\]
are situated in this ambient direct product, for arbitrary \(A,B\) and fixed nontrivial \(C\). The maps \(\iota_\alpha\) are placeholders for the cited construction, not specified embeddings supplied by this draft. The record neither gives that construction nor defines the intended scope of ``situated.'' Consequently, an exact characterization target cannot be recovered from the supplied record alone; the embedding construction, wreath-product convention, and intended subgroup relationships require clarification from the cited source.
\subsection{Short English statement}
For a fixed nontrivial group \(C\), Merzlyakov's problem asks how the split extensions associated with arbitrary groups \(A\) and \(B\) occur inside \(A\times\operatorname{Aut}(B\wr C)\) under a cited common embedding construction. The supplied record does not specify the construction, the wreath-product convention, or precisely which relationships among the embedded subgroups should be described.
\subsection{Sources}
\begin{itemize}
\item Local UnsolvedMath registry, supplied complete source record: permanent ID \(700248\), exact source problem code \texttt{catalogue entry}, title ``Kourovka Notebook Problem 4.40.'' The registry labels the problem open; this is supplied metadata, not an independently verified mathematical assessment.
\item Yu. I. Merzlyakov, proposer, 1973. The supplied record identifies the original source as \emph{The Kourovka Notebook}, 4th issue (1973), Problem 4.40, and its later presentation in the 21st issue edition (2026), arXiv:1401.0300v48. The collection metadata names E. I. Khukhro and V. D. Mazurov as editors.
\item Yu. I. Merzlyakov, \emph{Algebra and Logic}, volume 9, no. 5 (1970), pp. 326--337. The supplied statement cites this reference for the unified embedding construction; no article title or construction details are supplied, and the reference has not been independently checked here.
\item Supplied discussion: Achyuth Jayadevan, \emph{Preprint}, 2026, \url{https://zenodo.org/records/22727276}. The record explicitly labels the asserted description an ``unverified AI solution'' claim, not an established solution.
\item Supplied status review dated \texttt{2026-10-16}: canonical status \texttt{open}, evidence category \texttt{unrefereed\_claim}. Its note reports a \(\#\) marker for the unverified AI solution claim in the October 2026 21st issue, arXiv:1401.0300v48. This is supplied dated curation, not independently verified source or literature assessment.
\end{itemize}
\end{document}
